\chapter{Noncritical Probability Quantifiers and Almost-Everywhere Elimination}
\label{sec:p15-probability-elimination}

Keisler--Lotfallah's held author manuscript develops this construction on
pp.~8--28; p.~29 completes its references
\cite{KeislerLotfallah2009AuthorProbability}. The present account retains
its critical-value test, ordinary existential witnesses, parameter equality
patterns, and distinction between weak and strong almost-everywhere
reduction. The proofs below are independent reconstructions; the
concentration argument deliberately does not assume independence of
evaluations at overlapping tuples. This later finite-model paper does not
replace the still-unavailable 1985 chapter.

Use the two sample spaces and satisfaction clauses of
\cref{def:p15-pq-models,def:p15-pq-language}. For the zero--one statements
here the finite relational vocabulary has positive arities and no
constants or function symbols. This guard matters: a nullary relation
whose model probability is \(p\in(0,1)\) is already a quantifier-free
sentence of probability \(p\). The unqualified extension to nullary
symbols would therefore be false. Parameters \(p_R\) are initially fixed
strictly between zero and one.

\begin{hbdefinition}{def:p15-pe-types}{Equality diagrams, complete types, and conditional limits}
For a list \(\bar x=(x_1,\ldots,x_k)\) of distinct variable symbols, an
equivalence relation \(E\) on its positions determines the equality
diagram \(D_E(\bar x)\): coordinates in the same class are equal and
coordinates in different classes are unequal. Write \(D(\bar x)\) for
the all-distinct diagram. A complete quantifier-free type \(\beta(\bar x)\)
specifies every relation atom on these variables and every equality,
consistently; it induces one such \(E\). The tuple \(\bar x_E\) replaces
each class by one variable, repeated at every position in that class.
Types and diagrams concern values, not merely variable spellings.

Extend the notation \(\mu_n(\theta(\bar x))\) by sampling both the random
structure and an independent uniform parameter tuple. For a
quantifier-free \(\alpha(\bar x,\bar y)\), with disjoint variable lists,
define
\[
 q_\beta(\alpha;\bar y)
   =\lim_{n\to\infty}
       \frac{\mu_n(\alpha(\bar x,\bar y)\land\beta(\bar x))}
            {\mu_n(\beta(\bar x))}.
\]
The denominator is positive once \(n\) realizes the equality classes.
The \(\bar y\)-critical values are the finitely many \(q_\beta\), over
all complete parameter types. The next result constructs the limit;
its definition does not assume that construction.
\end{hbdefinition}

\begin{hbtheorem}{thm:p15-pe-critical}{proved=annex:p15.pe.critical}{The actual critical-value computation}
Collapse the parameter equalities, evaluate parameter-only atoms by
\(\beta\), and evaluate equalities on otherwise distinct fresh coordinates.
Express the resulting Boolean formula in complete, mutually exclusive
truth patterns of its distinct remaining relation atoms. Then
\[
 q_\beta=\sum_{\text{satisfying patterns }v}
             \prod_{\text{distinct atom keys }t}
                  p_{R_t}^{v_t}(1-p_{R_t})^{1-v_t}.
\]
Identical atom keys are counted once. The empty Boolean cases give zero
or one. This finite sum proves existence of the conditional limit and
is a polynomial with integer coefficients in the atomic probabilities.
Each collapsed consistent type has a positive limiting probability.
For rational atomic probabilities the critical values are rational.
\end{hbtheorem}

\begin{hbdefinition}{def:p15-pe-content}{Noncritical syntax and quantifier-free content}
Noncriticality and \(\varphi^0\) are defined together, as in source
Definition~3.6. Quantifier-free formulas are noncritical and retain
themselves. Negation and conjunction are noncritical exactly when their
constituents are, and content commutes with these connectives. For an
ordinary single-variable binder, if \(\psi\) is noncritical, put
\[
 (\exists y\,\psi(\bar x,y))^0
   =\bigvee_{q_\beta(\psi^0;y)>0}\beta(\bar x)
       \ \vee\ \bigvee_{i=1}^{k}\psi^0(\bar x,x_i).
\]
For a probability binder its scope must be noncritical and its threshold
\(r\) must differ from every \(\bar y\)-critical value of \(\psi^0\);
then put
\[
 ((\exists^{\ge r}\bar y)\psi(\bar x,\bar y))^0
       =\bigvee_{q_\beta(\psi^0;\bar y)>r}\beta(\bar x).
\]
Empty disjunctions are false. Substitution is capture-avoiding.
For finitely many free variables, even an infinitary quantifier-free
Boolean combination can be replaced by a finite truth table on the
finitely many atoms. These clauses do not admit a critical tie by
changing \(>\) to \(\ge\).
\end{hbdefinition}

\begin{hbtheorem}{thm:p15-pe-uniform}{proved=annex:p15.pe.uniform}{Uniform elimination of one quantifier-free probability binder}
Let \(\alpha(\bar x,\bar y)\) be quantifier-free and let \(r\in(0,1)\)
be noncritical. With probability tending to one, the random structure
satisfies, simultaneously for every \(\bar a\in n^k\),
\[
 (\exists^{\ge r}\bar y)\alpha(\bar a,\bar y)
 \quad\longleftrightarrow\quad
       \bigvee_{q_\beta(\alpha;\bar y)>r}\beta(\bar a).
\]
All parameter equality patterns and all tuple collisions are included.
For a fixed formula the exceptional probability is at most a constant
times \(n^k e^{-c n}\), for some \(c>0\) and all sufficiently large \(n\).
The annex proves a finite-product bounded-difference estimate, including
its exponential-moment argument, rather than treating the tuple truth
indicators as independent.
\end{hbtheorem}

\begin{hbtheorem}{thm:p15-pe-elimination}{proved=annex:p15.pe.elimination}{Finitary noncritical elimination and its effective scope}
Every finitary noncritical \(\varphi(\bar x)\) satisfies
\[
 \lim_{n\to\infty}\mu_n
       \bigl(\forall\bar x(\varphi(\bar x)\leftrightarrow
                                      \varphi^0(\bar x))\bigr)=1.
\]
The ordinary existential clause is proved by actual fresh witnesses
and the explicit parameter-witness disjunction above. Structural
induction handles nested probability binders without exchanging them.
For sentences the content is true or false, giving a zero--one law.
For a fixed rational atomic-probability vector and rational threshold
syntax, noncriticality and content have finite recursive computations
(indeed primitive-recursive ones under usual finite syntax encodings).
No effectiveness for arbitrary real inputs is asserted.
\end{hbtheorem}

For example, for a binary relation with probability \(p\), the fresh
limit of \(R(x,y)\) is \(p\). A threshold below \(p\) has content true;
one above \(p\) has content false; equality is excluded. For the unary
scope \(y=x\), the fresh limit is zero, yet \(\exists y(y=x)\) has
content true through its parameter witness. This is why the second
disjunction in the ordinary existential clause cannot be dropped.

\begin{hbdefinition}{def:p15-pe-ae}{Weak and common-class reduction}
Two formulas are almost-everywhere equivalent when the model probability
of their universally closed equivalence tends to one. Weak reduction
allows a different asymptotic class for each formula. Strong reduction
requires one class \(C\) of finite structures, with
\(\mu_n(C\cap M_n)\to1\), on which every formula in the declared fragment
equals its chosen first-order content at every assignment.
Neither is a logical equivalence over all finite structures.
\end{hbdefinition}

\begin{hbtheorem}{thm:p15-pe-strong}{proved=annex:p15.pe.strong}{A common class for finitely many variables and thresholds}
Fix a variable budget \(k\) and a finite \(P_0\subset(0,1)\).
The noncritical fragment using these variables and only thresholds in
\(P_0\) strongly reduces to the \(k\)-variable quantifier-free fragment.
This includes well-founded infinitary conjunctions and disjunctions,
not just finitary formulas. The union over finite budgets and finite
threshold sets has formula-by-formula reduction and a zero--one law
for sentences. The proof constructs a finite set of primitive
quantifier-elimination identities before choosing the common class;
it does not intersect arbitrarily many asymptotic measure-one events.
\end{hbtheorem}

There is no common class for all thresholds, even with one variable:
for an unbiased unary \(R\), put \(r_m=1/2-1/m\), \(m\ge3\).
Every \((\exists^{\ge r_m}x)R(x)\) has true content, but their
conjunction is exactly the critical half-threshold assertion, of limiting
probability \(1/2\) by \cref{thm:p15-pq-limit-failures}.
Consequently an asymptotic measure-one class cannot realize all these
contents. This also supplies a noncritical infinitary counterexample
when infinitely many thresholds are allowed. The held manuscript's
p.~25 instead prints these below-half formulas as having asymptotic
measure zero and uses a disjunction; those two statements are not
adopted. The conjunction above is the independently checked correction.

\begin{hbdefinition}{def:p15-pe-parametric}{Parametric constraints and their actual conditioned model law}
A parametric sentence \(\pi\) is a finite conjunction of
\(\forall\bar z[D(\bar z)\to\tau(\bar z)]\), where \(\tau\) is a
Boolean combination of literals and every variable of \(\bar z\)
occurs in every literal. Assume \(\pi\) has models of every positive
finite size. The model measure is
\[
  (\mu_n\mid\pi)(B)=
            \frac{\mu_n(B\cap\{\mathcal A:\mathcal A\models\pi\})}
                 {\mu_n(\pi)}.
\]
Use only complete parameter types consistent with \(\pi\); their
conditional scope limits define \(q_\beta^\pi\), and the preceding
content clauses with those limits define \(\varphi^\pi\).
Conditioning does not preserve independence of individual atoms.
\end{hbdefinition}

\begin{hbtheorem}{thm:p15-pe-parametric}{proved=annex:p15.pe.parametric}{Parametric elimination through independent support blocks}
For this law, every finitary \(\pi\)-noncritical formula is
almost-everywhere equivalent to \(\varphi^\pi\), and its sentences have
a zero--one law. The limits \(q_\beta^\pi\) exist and are positive-type
conditional finite sums. The proof groups atoms by the exact set of
elements in their arguments: distinct support blocks remain independent
after conditioning, whereas atoms within one block can be correlated.
It constructs their normalized laws, the limits, the witness argument,
and concentration. It is not an inference from an unproved transfer of
the unconditioned theorem. The same finite primitive-identity argument
also gives the fixed-variable, finite-threshold common-class version.
\end{hbtheorem}

\begin{hbtheorem}{thm:p15-pe-model-boundaries}{proved=annex:p15.pe.model-boundaries}{What conditioning and the source's graph examples actually permit}
For a binary relation whose ordered ground atoms have probability \(p\),
conditioning on symmetry and irreflexivity makes the probability of a
distinct undirected edge
\[
             \frac{p^2}{p^2+(1-p)^2},
\]
not generally \(p\). This is an actual parametric class.
In contrast, in the ordinary vocabulary of one uncolored binary graph
relation, the class of all finite bipartite graphs cannot be a parametric
class in Definition~5.1's sense. The source's p.~22 mention of bipartite
graphs is therefore not adopted under this uncolored interpretation;
an additional presentation would have to be specified and justified.
If equality literals are permitted, a model whose size equals the
largest relation arity need not guarantee models of every size:
in a unary vocabulary the parametric constraint
\(\forall x y[D(x,y)\to x=y]\) has only singleton models.
The all-size hypothesis in \cref{def:p15-pe-parametric} is therefore
retained explicitly rather than inferred from that arity-only criterion.
The nullary-symbol extension of the zero--one assertions is also false,
by the one-atom countermodel already identified above.
\end{hbtheorem}

\begin{hbdefinition}{def:p15-pe-sequences}{The source's size-dependent noncriticality}
Now allow \(p_R(n)\in(0,1)\) and a threshold sequence \(r(n)\in(0,1)\).
For each complete \(\beta\), write
\(q_{\beta,n}=\mu_n[\alpha(\bar x,\bar y)\mid\beta(\bar x)]\);
these sequences need not converge. A sequence \(r(n)\) is critical for
\(\alpha\) if some \(\beta\) satisfies both
\[
 \liminf(q_{\beta,n}-r(n))\le0,\qquad
 \liminf(r(n)-q_{\beta,n})\le0.
\]
Otherwise the probability content selects the types with
\(\liminf(q_{\beta,n}-r(n))>0\).
For an ordinary existential with noncritical scope, every \(\beta\)
must satisfy either \(\liminf q_{\beta,n}>0\), or the logical validity
\[
 \forall\bar x\bigl(\beta(\bar x)\to
                       \bigvee_{i=1}^{k}\psi^0(\bar x,x_i)\bigr).
\]
Its content is the disjunction of the positive-liminf types and the
parameter-witness terms. These are source Definitions~7.1--7.2,
not merely the condition \(r(n)\ne q_{\beta,n}\) at each finite size.
For example, if a unary fresh-scope mean is the constant \(p\), then
a threshold \(p-1/n\) for all sufficiently large \(n\) is still critical:
both displayed liminf differences are zero, despite strict inequality
at every sufficiently large finite size.
\end{hbdefinition}

\begin{hbtheorem}{thm:p15-pe-sequences}{proved=annex:p15.pe.sequences}{Size-dependent elimination at the stated separation hypotheses}
With those clauses and a fixed finite variable budget and finite set of
threshold sequences, the noncritical infinitary fragment has a common
asymptotic class and equals its quantifier-free content there.
The union of these fragments has weak reduction and a zero--one law.
No global lower bound on all \(p_R(n)\) is required for this selected
result: the existential and threshold clauses supply the necessary
eventual margins themselves. The independent ground-atom law is
retained, and parametric conditioning is not silently added here.
\end{hbtheorem}

The numerical threshold interpretation is exact:
\((\exists^{\ge r(n)}\bar y)\alpha\) counts at least
\(\lceil r(n)n^{|\bar y|}\rceil\) tuples. For a single variable,
the one-witness threshold \(1/n\) agrees semantically with ordinary
existence at \(n>1\); the source's open interval excludes its value at
\(n=1\). Their noncriticality tests still differ. No elimination claim
is made at omitted critical ties, for arbitrary unbounded infinitary
threshold inventories, or for the full unrestricted language.
Nor do these asymptotic constructions provide a completeness calculus
or the general measure/conditional-expectation treatment of the
original 1985 allocation.
